\label{sec:estrategia}

\subsection{Conversão, perdas e ordenação temporal}

A variável dependente é a taxa anual de conversão bruta da vegetação nativa para agricultura. Para o município $i$ e o ano $t$, divide-se a área dessa transição pelo estoque nativo em $t-1$, multiplicando-se por cem. A unidade é ponto percentual (p.p.) do estoque inicial. A mesma construção define o desfecho complementar de pastagem. Ao observar diretamente as transições, a medida evita que regeneração em outras parcelas compense a abertura do ano, como ocorreria numa diferença líquida entre estoques. O destino registrado é o uso imediatamente posterior à vegetação; a classificação não informa a cultura plantada, a legalidade da abertura ou os agentes que a realizaram.

A quebra é construída separadamente para cada cultura $c$, comparando o desempenho realizado com uma referência histórica municipal. Utiliza-se $r^c_{is}=Q^c_{is}/AP^c_{is}$, em que $Q$ é a produção e $AP$, a área plantada na safra anual $s$. Quando as áreas são positivas, essa razão pode ser escrita como rendimento por hectare colhido multiplicado pela proporção colhida da área plantada. Essa decomposição dá sentido à escolha: um município pode apresentar perda mesmo que o rendimento das parcelas que chegaram à colheita se mantenha. A medida também pode variar com manejo, composição de áreas ou erro de declaração, de modo que seu déficit não equivale automaticamente a dano climático.

A referência $\widehat r^c_{is}$ resulta de uma tendência linear estimada com os dez anos anteriores a $s$, exigindo pelo menos sete observações válidas e previsão positiva. A tendência admite evolução gradual do desempenho, ao passo que a janela móvel limita a influência de um passado produtivo distante. Nenhuma observação futura participa da previsão. Descontinuidades territoriais identificadas reiniciam o histórico, evitando que a mudança da unidade seja interpretada como perda. A intensidade é definida por:
\begin{equation}
L^c_{is}=\max\left\{0,\,1-\frac{r^c_{is}}{\widehat r^c_{is}}\right\}.
\label{eq:perda}
\end{equation}
Assim, $L=0,20$ indica desempenho 20\% abaixo da referência local, e realizações acima dela recebem zero na especificação principal. Ausência de cultivo e informação insuficiente permanecem sem perda calculada. Uma versão binária identifica déficits de pelo menos 20\%; alternativas utilizam tendências de cinco ou quinze anos, com mínimos de quatro ou onze observações, mediana histórica de dez anos com mínimo de sete, área colhida e limiares de 10\% ou 30\%. A referência é uma estimativa, e não uma produtividade que teria sido observada com certeza na ausência de choque.

O alinhamento principal associa a transição dos mapas entre $t-1$ e $t$ à perda da safra $t-2$ e ao preço médio de $t-1$. Para a conversão de 2001, por exemplo, utiliza-se a perda de 1999, cuja referência pode mobilizar 1989--1998, e o preço de 2000. Para 2024, utilizam-se perda de 2022 e preço de 2023. A separação reduz a sobreposição entre o desempenho medido e a abertura que se pretende descrever, mas os dados anuais não permitem ordenar todas as decisões dentro de cada ano. A PAM agrega as safras de milho, e os mapas não datam cada supressão. A alternativa com perda em $t-1$ é, por isso, uma verificação distinta, sujeita a maior sobreposição temporal.

Define-se $p^c_{t-1}$ como o logaritmo do preço real, centrado na média dos 24 valores anuais usados no painel. A capacidade $Z^c_i$ é padronizada pela média e pelo desvio-padrão dos 141 municípios, conservados em todas as amostras. Essa padronização permite comparar coeficientes sem que a unidade de capacidade mude quando municípios são excluídos. Uma variação de 10\% no preço corresponde a $\ln(1,10)$ na escala logarítmica. O modelo em nível do logaritmo relaciona o patamar real do incentivo à conversão; a variação anual e a razão entre preços respondem a perguntas diferentes e entram apenas como alternativas identificadas separadamente.

\subsection{Contrastes de capacidade e perdas}

O núcleo procura distinguir duas formas de heterogeneidade. A primeira compara a associação com preços entre municípios de diferentes capacidades potenciais; a segunda compara essa associação diante de diferentes perdas anteriores. Para cada cultura, estima-se:
\begin{equation}
\begin{split}
y^{ag}_{it}={}&\alpha_i+\lambda_{b(i),t}
 +\beta_c p^c_{t-1}Z^c_i
 +\theta_c L^c_{i,t-2}
 +\delta_c p^c_{t-1}L^c_{i,t-2}\\
 &+\gamma_c Z^c_i\tau_t+\rho_c R_{it}+\varepsilon^c_{it},
\end{split}
\label{eq:nucleo}
\end{equation}
em que $\alpha_i$ são efeitos municipais, $\lambda_{b(i),t}$ são efeitos do bioma predominante por ano, $\tau_t=(t-2001)/10$ e $R_{it}$ é a anomalia anual de precipitação, em centenas de milímetros. Os efeitos municipais absorvem características locais constantes; os de bioma por ano absorvem choques comuns aos municípios de cada grupo no mesmo ano. A tendência diferenciada por capacidade permite que municípios de maior potencial sigam trajetória gradual distinta, além da associação com o preço. A chuva representa uma condição contemporânea que pode afetar diretamente vegetação e conversão, sem ser utilizada como instrumento para a perda.

Os coeficientes $\beta_c$ e $\delta_c$ são comuns aos biomas no núcleo. Os efeitos de bioma por ano condicionam a comparação, mas não correspondem a inclinações próprias para Amazônia, Cerrado e Pantanal. As extensões territoriais descritas adiante examinam outra dimensão, relacionada às frações de exposição e aos marcos institucionais, com os limites de suporte de cada cultura.

Os níveis isolados de preço e capacidade não aparecem porque são absorvidos pelos efeitos fixos. Sua interação continua identificável se houver variação suficiente de capacidade entre municípios e de preços ao longo do tempo, após os demais termos. A perda varia nas duas dimensões e permanece na equação junto com preço×perda. O parâmetro $\theta_c$ descreve sua associação no preço centrado igual a zero; H1b se refere a $\delta_c$, que altera a inclinação em relação ao preço. A presença de uma associação entre perda e conversão, sem essa interação, não seria um teste da sub-hipótese.

Para H1a, o contraste entre municípios separados por um desvio-padrão de capacidade, diante de uma alta de 10\% no preço, é $\beta_c\ln(1,10)$ p.p. A hipótese prevê $\beta_c>0$. Para H1b, a diferença associada a dez pontos percentuais adicionais de perda é $\delta_c\ln(1,10)\times0,10$ p.p.; seu sinal é uma questão empírica. Esses contrastes não recuperam o efeito médio de uma alta de preços em Mato Grosso, absorvido pelos efeitos temporais. Eles descrevem diferenças condicionais de associação, com as demais variáveis da equação mantidas constantes.

A especificação B corresponde à Equação~\ref{eq:nucleo}. A especificação A exclui os dois termos de perda e é estimada tanto na amostra disponível originalmente quanto nas mesmas chaves de B. A comparação entre A original e A comum identifica a alteração associada ao recorte, enquanto A comum e B mostram a alteração decorrente da inclusão dos termos, conservando as observações. A amostra elegível pela perda contém 1.488 pares na soja e 2.381 no milho. A remoção de grupos unitários sem contribuição à estimação com esses efeitos fixos deixa, respectivamente, 1.459 pares em 106 municípios e 2.376 em 134. A original reúne 2.767 pares em 134 municípios. A versão com limiar de 20\% utiliza as mesmas chaves de B.

Essa distinção é relevante para o alcance territorial. No Pantanal, B dispõe de oito pares em dois municípios para soja e cinquenta em três para milho. O painel não fornece base comparável à dos demais biomas para interpretar uma resposta própria dessa região. O histórico também seleciona localidades com cultivo e registro suficientemente contínuos. A análise das perdas diz respeito a esse conjunto, sem representar integralmente municípios de entrada recente na agricultura ou unidades sem informação válida.

\subsection{Inferência e verificações de estabilidade}

A multiplicação dos preços anuais por características municipais amplia a variação das interações, mas não cria novos choques independentes de preço. Apresentam-se três procedimentos de inferência: agrupamento por município, agrupamento simultâneo por município e ano e covariância robusta à heterocedasticidade e à autocorrelação (HAC), com banda dois sobre os escores anuais da regressão completa residualizada. Neste último, a aplicação de Frisch--Waugh--Lovell preserva os coeficientes de B; utilizam-se pesos de Bartlett, sem pré-branqueamento ou ajuste adicional de pequena amostra. Os graus de liberdade são, respectivamente, $G-1$, $\min(G,T)-1$ e $T-1$, em que $G$ é o número de municípios e $T$, o de anos. Com 24 anos, as duas últimas referências utilizam 23 graus de liberdade.

O agrupamento municipal admite dependência serial dentro da unidade; o agrupamento duplo acrescenta dependência entre municípios no mesmo ano. O HAC dos escores anuais examina a persistência temporal da contribuição agregada à estimação. Nenhum procedimento garante acomodar toda forma de dependência espacial e temporal, especialmente com poucos anos e interações concentradas. Por isso, a apresentação conserva a divergência entre as inferências. Intervalos são de 95\%, e ajustes de Holm são aplicados a famílias explicitadas: os dois gradientes H1a de B, separadamente dos dois modificadores H1b. Nas extensões, as famílias acompanham o conjunto de contrastes efetivamente examinado; os intervalos apresentados continuam individuais.

As verificações incluem dezesseis definições por cultura, com variações da referência produtiva, intensidade ou ocorrência da perda, calendário, clima, efeitos temporais, tendência por capacidade, janela e preço. Comparações intermediárias nas mesmas chaves separam, quando necessário, modificação da fórmula de alteração de amostra. Também se examinam exclusões locais e anuais, reamostragem restrita por ano e perturbações da referência histórica. Essas últimas verificam a sensibilidade ao indicador gerado, sem converter sua distribuição em intervalo que incorpore toda a incerteza do desenho.

Como antecedentes de H1a, retomam-se exercícios sem perdas para soja que substituem o potencial de alto insumo pelo rendimento potencial em baixo insumo ou pelo índice de adequação GAEZ em alto insumo. As medidas são padronizadas e interagem com preço e tendência, mantendo conversão agrícola, chuva, efeitos municipais e de bioma por ano, em 2.767 pares. Outro exercício usa a média de 1985--2000 da proporção municipal ocupada por agricultura, calculada com as classes 20, 39, 41 e 62 do MapBiomas, com os mesmos tipos de interação e ajustes. Esses modelos descrevem sensibilidade à medida de exposição, com erros-padrão agrupados por município. Não incorporam a perda nem substituem a comparação de A e B na amostra comum.

Três períodos --- 2001--2005, 2006--2011 e 2012--2024 --- permitem verificar se uma associação agregada esconde trajetórias distintas. Interagem-se os termos de interesse com os períodos, preservando os constitutivos e os ajustes de tendência previstos na especificação. Os testes conjuntos examinam tanto a nulidade dos modificadores quanto sua igualdade entre períodos. São comparações exploratórias: os períodos iniciais contêm poucos preços e poucas perdas intensas, sobretudo na soja. Seu uso não transforma as datas em mudanças exógenas ou em avaliação de políticas.

\subsection{Destinos, incentivos conjuntos e instituições}

A extensão para pastagem conserva fórmula e observações de B. A diferença agricultura menos pastagem é estimada diretamente, de modo que sua variância incorpora a covariância entre os destinos. Outra extensão acrescenta a interação entre capacidade e perda e a interação tripla entre preço, capacidade e perda, permitindo que a modificação de H1b também varie com o potencial. Os dois termos entram juntos para respeitar a hierarquia da interação. Os contrastes condicionais são avaliados na distribuição observada; percentis marginais não asseguram observações em todas as combinações de capacidade e perda.

O preço relativo é $\ln[P^c/(1000P^{carne})]$, com o fator mil compatibilizando cotações por tonelada e por quilograma. Essa medida compara incentivos brutos entre produtos, sem descontar custos. Testa-se a restrição de coeficientes opostos que ela impõe a uma versão com preços separados. Estima-se ainda uma versão conjunta de soja e milho na interseção de 1.450 pares em 105 municípios, mantendo referências individuais nessas mesmas chaves. A correlação de 0,948 entre as séries de preços justifica cautela ao separar sua contribuição, mesmo quando a matriz de interações mantém posto completo.

A exposição territorial à Moratória utiliza a fração municipal no bioma Amazônia interagida com o período desde 2007, com alternativa desde 2006. O pacto foi anunciado em julho de 2006; a alteração do marco de desmatamento para julho de 2008 ocorreu em dezembro de 2014. Nenhuma dessas datas equivale a um início uniforme de cumprimento municipal. A medida representa alcance potencial do compromisso sobre a soja, sem observar fornecedores e transações; para milho, caracteriza heterogeneidade no território alcançado pelo pacto de outra cultura.

Para a reforma do Código Florestal, utiliza-se a fração de formações florestais de 2000 na área municipal total, inclusive municípios com valor zero, combinada ao período desde 2013 e, alternativamente, desde 2012. A medida reúne classes 3 e 6 do MapBiomas e não corresponde a percentual jurídico de Reserva Legal, passivo ou situação cadastral. As extensões institucionais acrescentam, separadamente, as interações de preço e capacidade, ou de preço e perda, com exposição e período posterior, com todos os termos de ordem inferior que não são absorvidos e tendência pela exposição. Sua interpretação depende da variação residual disponível e da possibilidade de outras mudanças acompanharem a exposição e a data.

O desenho permanece associativo. Preços internacionais não medem o retorno líquido recebido localmente; clima, infraestrutura, tecnologia e composição produtiva podem variar junto com as interações. A defasagem da perda ordena medidas, mas não elimina choques persistentes ou decisões anteriores que afetem simultaneamente desempenho e conversão. Controles como crédito, renda e produção posteriores à quebra poderiam representar consequências do próprio evento, alterando o contraste. Fiscalização observada também pode responder à pressão de abertura. Esses limites orientam a escolha do núcleo parcimonioso e a interpretação das extensões, sem supor que acrescentar controles produza automaticamente identificação causal.
